\subsubsection{Hamiltonian Recap}\label{sec:hamiltonian-recap}

Before connecting QUBO problems to quantum systems, we need to introduce one of the most important concepts in quantum mechanics: the Hamiltonian. Very informally, the \definition{Hamiltonian} is the \textbf{mathematical object} that \hl{describes the energy of a quantum system and determines how the system evolves over time.}

\highspace
\begin{flushleft}
    \textcolor{Green3}{\faIcon{question-circle} \textbf{What is a Hamiltonian?}}
\end{flushleft}
First of all, the Hamiltonian is an \textbf{operator}. An operator is simply a mathematical object that acts on another object and transforms it in some way.

\highspace
For example, the derivative is an operator:
\begin{equation*}
    D\left[f(x)\right] = \frac{d}{dx} f(x)
\end{equation*}
Here, the operator $D$ takes a function $f(x)$ and produces its derivative $\frac{d}{dx} f(x)$. For instance, if:
\begin{equation*}
    f(x) = x^2 \; \Rightarrow \; D\left[f(x)\right] = 2x
\end{equation*}
Similarly, a Hamiltonian takes a quantum state and acts on it. Imagine a physical system, such as a collection of interacting qubits. At any instant, the system has some energy. The Hamiltonian is the mathematical object that tells us:
\begin{enumerate}
    \item \textbf{How much energy every possible state has}
    \item \textbf{How the system evolves over time}
\end{enumerate}
So, the Hamiltonian is the \emph{rulebook} that \hl{assigns an energy to every possible state of a system.}

\highspace
The landscape introduced for QUBO (\autopageref{fig:energy-landscape-xor}) is a great \textbf{analogy}. For XOR, we had an energy landscape where each point corresponds to a particular assignment of the variables $x_1$ and $x_2$. Each state (point) has an associated energy. The Hamiltonian does exactly the same thing for a quantum system. In some sense, the Hamiltonian \textbf{is the mathematical representation of the energy landscape}.

\highspace
\begin{flushleft}
    \textcolor{Green3}{\faIcon{tools} \textbf{Quantum Mechanics View}}
\end{flushleft}
Historically, in \textbf{classical mechanics}, the Hamiltonian is often written as:
\begin{equation*}
    H = T + V
\end{equation*}
Where $T$ is the kinetic energy and $V$ is the potential energy. For example, for a ball:
\begin{equation*}
    H = \frac{1}{2} mv^{2} + mgh
\end{equation*}
It represents the total energy of the system, which is the sum of kinetic and potential energy, called Hamiltonian.

\highspace
However, in \textbf{quantum mechanics} things become more complicated. A quantum state is represented by a vector $\ket{\psi}$. The Hamiltonian becomes a \textbf{matrix} (more generally, an \textbf{operator}): $\hat{H}$. Instead of assigning a single number directly, it \textbf{acts on quantum states} $\hat{H}\ket{\psi}$. The result tells us information about the \textbf{energy of the state} and \textbf{how} the state \textbf{evolves over time}. Unlike quantum gates, the Hamiltonian describes the physical system, exists naturally within it, and determines its energy and evolution. It does not perform a computation or change the state in a specific way.

\highspace
\begin{flushleft}
    \textcolor{Green3}{\faIcon{hat-wizard} \textbf{Why does the Hamiltonian determine time evolution?}}
\end{flushleft}
We have seen that the Hamiltonian assigns an energy to every state naturally. But it also \textbf{determines how the system evolves over time}. Let's briefly explain why.

\highspace
We already know that a quantum state is \ket{\psi}. Suppose at time $t=0$ we have a state \ket{\psi(0)}. Now, the question: \emph{what will the state be one second later?} In classical mechanics we would use Newton's laws to determine the future state, like $F = ma$, which tells us how the position and velocity of an object change over time. Instead, in quantum mechanics, we use the Schrödinger equation:
\begin{equation}
    i\hbar \frac{\partial}{\partial t} \ket{\psi(t)} = \hat{H} \ket{\psi(t)}
\end{equation}
The important thing is not the formula itself (which we will not analyze in detail in this course), but the fact that \textbf{Schrödinger equation predicts future quantum state}. And the \hl{object appearing inside the Schrödinger equation is the Hamiltonian}. Therefore, if we know the Hamiltonian, we can predict the future evolution of the quantum system.

\highspace
In other words, \hl{if we know the Hamiltonian} $\hat{H}$ of a quantum system, then the \hl{Schrödinger} equation \hl{tells us how the quantum state evolves} in time. Therefore, knowing the \hl{Hamiltonian allows us to predict the future} state of the system.

\highspace
\begin{flushleft}
    \textcolor{Green3}{\faIcon{question-circle} \textbf{But how does a Hamiltonian store energies?}}
\end{flushleft}
First of all, let's remember where we came from. We built an energy landscape (\autopageref{sec:qubo-as-an-energy-minimization-problem}) where each state (configuration of variables) has an associated energy. The optimization problem is:
\begin{equation*}
    \min E(x)
\end{equation*}
Where we want to find the state with the lowest energy.

\highspace
\textcolor{Green3}{\faIcon{question-circle} \textbf{Now, what do we need from a quantum system?}} If we want a quantum system to solve this problem, we need a way to tell the quantum system: ``\emph{Hey, for this state, the energy is $E_1$. For that state, the energy is $E_2$. For that other state, the energy is $E_3$}''. In other words, we need a way to \textbf{encode the energy landscape into the quantum system}. So, \emph{how do we store these energies inside a quantum system?} Simple, we use the Hamiltonian.

\highspace
\textcolor{Green3}{\faIcon{question-circle} \textbf{But how does a Hamiltonian store energies?}} This is where eigenvalues and eigenvectors come into play. When we have a Hamiltonian $\hat{H}$, we can find its eigenvalues and eigenvectors:
\begin{equation}\label{eq:hamiltonian-eigenvalues-eigenvectors}
    \hat{H} \ket{\phi_i} = E_i \ket{\phi_i}
\end{equation}
First of all, we apply the Hamiltonian to the eigenvector $\ket{\phi_i}$. The result is the same vector $\ket{\phi_i}$ multiplied by a number $E_i$. This number $E_i$ is called the \textbf{eigenvalue} corresponding to the \textbf{eigenvector} $\ket{\phi_i}$.

\highspace
For example, if we have a Hamiltonian $\hat{H}$:
\begin{equation*}
    \hat{H} = \begin{bmatrix}
        1 & 0 \\
        0 & 3
    \end{bmatrix}
\end{equation*}
And the state:
\begin{equation*}
    \ket{0} = \begin{bmatrix}
        1 \\
        0
    \end{bmatrix}
\end{equation*}
Applying the Hamiltonian to the state gives us:
\begin{equation*}
    \hat{H} \ket{0} = \begin{bmatrix}
        1 & 0 \\
        0 & 3
    \end{bmatrix} \begin{bmatrix}
        1 \\
        0
    \end{bmatrix} = \begin{bmatrix}
        1 \\
        0
    \end{bmatrix} = 1 \ket{0}
\end{equation*}
Here, the eigenvalue is $E_0 = 1$ corresponding to the eigenvector $\ket{0}$. Similarly, for the state:
\begin{equation*}
    \ket{1} = \begin{bmatrix}
        0 \\
        1
    \end{bmatrix}
\end{equation*}
Applying the Hamiltonian gives us:
\begin{equation*}
    \hat{H} \ket{1} = \begin{bmatrix}
        1 & 0 \\
        0 & 3
    \end{bmatrix} \begin{bmatrix}
        0 \\
        1
    \end{bmatrix} = \begin{bmatrix}
        0 \\
        3
    \end{bmatrix} = 3 \ket{1}
\end{equation*}
Here, the eigenvalue is $E_1 = 3$ corresponding to the eigenvector $\ket{1}$. It means that the state $\ket{0}$ has energy $1$ and the state $\ket{1}$ has energy $3$.

\highspace
The key point is that \textbf{the energy of a quantum system is encoded in the eigenvalues of the Hamiltonian}. Indeed, if \ket{\psi_{i}} is an eigenvector of the Hamiltonian, then the corresponding eigenvalue $E_i$ is the energy of that state.

\highspace
\begin{flushleft}
    \textcolor{Red2}{\faIcon{exclamation-triangle} \textbf{Hamiltonian must be Hermitian}}
\end{flushleft}
First of all, a matrix is \textbf{Hermitian} if it is \textbf{equal to its conjugate transpose}:
\begin{equation*}
    H = H^{\dagger}
\end{equation*}
Where $\dagger$ means the transpose and complex conjugate of the matrix. For example:
\begin{equation*}
    A = \begin{bmatrix}
        1 & 2 \\
        2 & 3
    \end{bmatrix} = A^{\dagger}
\end{equation*}
Similary, the complex matrix:
\begin{equation*}
    B = \begin{bmatrix}
        1 & i \\
        -i & 2
    \end{bmatrix}
    \to
    B^{T} = \begin{bmatrix}
        1 & -i \\
        i & 2
    \end{bmatrix}
    \to
    B^{\dagger} = \begin{bmatrix}
        1 & i \\
        -i & 2
    \end{bmatrix} = B
\end{equation*}
An interesting property of Hermitian matrices is that they have \textbf{real eigenvalues}. Taking the previous example, the eigenvalues of $A$ are $2 \pm \sqrt{5}$, which are real numbers. The eigenvalues of $B$ are $\approx 2.61$ and $\approx 0.38$, which are also real numbers. Also, the \textbf{eigenvectors corresponding to different eigenvalues are orthogonal} (\autopageref{eq:inner-product}).

\highspace
\textcolor{Green3}{\faIcon{question-circle} \textbf{Why must energy be real?}} Suppose we measure the energy of a particle. The experimental result will always be a real number. Now imagine the result of the measurement returned a complex number, like $3 + 7i$ joules. \emph{What would that even mean physically?} \emph{Is the particle more energetic because of the $7i$?} \emph{Can we extract $7i$ Joules of work?} There is no physical interpretation for a complex energy. Therefore, \textbf{measurable quantities must be real}, and energy is a measurable quantity, $E_{i} \in \mathbb{R}$. This implies that the \hl{Hamiltonian must have real eigenvalues}. In quantum mechanics, this is guaranteed by requiring that the Hamiltonian is a \textbf{Hermitian operator}.

\highspace
\textcolor{Green3}{\faIcon{question-circle} \textbf{Why is it useful that the energy eigenstates of the Hamiltonian are orthogonal?}} Let's \textbf{suppose} energy eigenstates were \textbf{not orthogonal}. Imagine we have two energy eigenstates $\ket{\phi_1}$ and $\ket{\phi_2}$ with two different energies. If they were not orthogonal, then projecting onto one state would partially affect the other. Then, the decomposition:
\begin{equation*}
    \ket{\psi} = c_1 \ket{\phi_1} + c_2 \ket{\phi_2}
\end{equation*}
Would become \emph{messy} (i.e., the coefficients $c_1$ and $c_2$ would not be independent). The coefficients would not be easy to extract, and the energy of the state would not be well-defined.

\highspace
Instead, \textbf{orthogonality gives clean probabilities}. Suppose:
\begin{equation*}
    \ket{\psi} = c_1 \ket{\phi_1} + c_2 \ket{\phi_2}
\end{equation*}
Because the eigenstates are orthonormal:
\begin{equation*}
    \braket{\psi_{0}}{\psi_{1}} = 0
\end{equation*}
We can \emph{easily} extract the coefficients:
\begin{equation*}
    P\left(E_0\right) = \left|c_1\right|^2, \quad P\left(E_1\right) = \left|c_2\right|^2
\end{equation*}
No cross terms appear. Precisely, \textbf{each energy eigenstate evolves independently}. Again, because the eigenstates are orthogonal, the evolution of one state does not affect the other.

\highspace
In simple terms, if our Hamiltonian has two energies $E_{1} \neq E_{2}$, when we measure energy, we want the outcomes to be either $E_{1}$ or $E_{2}$, with probabilities $P(E_1)$ and $P(E_2)$. If the eigenstates were not orthogonal, the measurement outcomes would not be well-defined, and we would not be able to assign clear probabilities to each energy level. Therefore, \textbf{orthogonality of energy eigenstates is crucial for a consistent physical interpretation of quantum mechanics}.